\documentclass[11pt]{article}

\usepackage[utf8]{inputenc}
\usepackage[T1]{fontenc}
\usepackage[spanish,es-tabla]{babel}
\usepackage[margin=2.2cm]{geometry}
\usepackage{amsmath}
\usepackage{amssymb}
\usepackage{booktabs}
\usepackage{longtable}
\usepackage{xcolor}
\usepackage{titlesec}
\usepackage{enumitem}
\usepackage{hyperref}
\hypersetup{colorlinks=true, linkcolor=blue!50!black, urlcolor=blue!50!black}

\definecolor{slidecolor}{RGB}{0,90,140}
\definecolor{timecolor}{RGB}{150,40,40}
\definecolor{cutcolor}{RGB}{160,110,0}

\titleformat{\section}{\normalfont\Large\bfseries\color{slidecolor}}{Slide \thesection}{0.6em}{}
\titlespacing*{\section}{0pt}{1.4em}{0.6em}

\newcommand{\tiempo}[1]{{\small\color{timecolor}\textbf{Tiempo sugerido: #1}}\par\vspace{0.3em}}
\newcommand{\recortable}{{\small\color{cutcolor}\textbf{[Recortable en vivo si falta tiempo]}}\par\vspace{0.3em}}
\newcommand{\respaldo}{{\small\color{cutcolor}\textbf{[Material de respaldo --- mostrar solo si preguntan]}}\par\vspace{0.3em}}

\title{Notas del expositor\\\large Ataques de envenenamiento de conocimiento a sistemas KG-RAG}
\author{Ingrid Alicia Yábar Geldres}
\date{Reproducción de Zhao et al. (2025), arXiv:2507.08862}

\begin{document}
\maketitle

\begin{abstract}
\noindent Este documento acompaña a \texttt{presentacion.tex} (23 diapositivas) con el guion
hablado de cada una, incluyendo el detalle completo de las ecuaciones 4--6 que se
fusionaron en una sola diapositiva. Los tiempos sugeridos suman
\textbf{\textasciitilde 20 minutos} de contenido, dentro de la ventana de 15--20 minutos
asignada; la diapositiva marcada como recortable puede omitirse si el tiempo aprieta. Las
dos diapositivas de apéndice no se presentan en el flujo normal: son material de respaldo
para preguntas puntuales.
\end{abstract}

\tableofcontents
\vspace{1em}
\hrule

% ============================================================
\section{Portada}
\tiempo{0:20}
Saludo breve, nombre, y una frase de encuadre: \emph{``Voy a presentar la reproducción de
un ataque de envenenamiento de conocimiento contra sistemas KG-RAG, del artículo de Zhao
et al. (2025).''} No hay contenido técnico en este slide; sirve solo para situar al
público antes de la diapositiva 3, que es la más importante de toda la exposición.

% ============================================================
\section{Contenido}
\tiempo{0:20}
Recorrer el índice en una frase por sección, sin detenerse: motivación, marco teórico,
método implementado, réplica del experimento, desviaciones y limitaciones, conclusiones.
No profundizar en nada todavía --- es solo un mapa.

% ============================================================
\section{Nota sobre la disponibilidad de código}
\tiempo{0:45}
\textbf{Esta es la diapositiva más importante de la exposición y debe decirse con
seguridad, no como disculpa.} El punto central: el artículo de Zhao et al. (2025) no
publica un repositorio de código ---se verificó en el PDF, en el HTML de arXiv y en la
página de resumen--- así que, a diferencia de otros artículos de la lista del curso, aquí
no es posible clonar y ejecutar el código original porque ese código no existe
públicamente. Por eso todo lo que sigue es una \textbf{implementación propia}, construida
directamente a partir de la descripción metodológica y de las ecuaciones 1 a 7 del
artículo, no a partir de un código fuente ajeno. Cerrar remitiendo a que las desviaciones
frente al método original están documentadas de forma explícita más adelante (diapositiva
18), para que la implementación sea auditable frente al texto original.

\emph{Si el profesor pregunta ``¿por qué no clonaste el repo?'' aquí mismo}: repetir que se
buscó activamente y no existe, y que se avisó de esto con anticipación (ver sección de
preguntas frecuentes al final de este documento).

% ============================================================
\section{¿Por qué importa la seguridad de KG-RAG?}
\tiempo{0:45}
Encadenar las cuatro viñetas como una idea continua, no como lista suelta: RAG dejar de
depender solo de lo memorizado en el entrenamiento y apoyarse en una fuente externa; cuando
esa fuente es un grafo de conocimiento (KG), el sistema puede razonar explícitamente,
salto a salto, entre entidades --- eso es KG-RAG, y es lo que lo hace interpretable. Pero
la misma edición del grafo que permite mantenerlo actualizado es la puerta de entrada del
atacante: una sola tripleta falsa, bien conectada, puede desviar una cadena de razonamiento
completa sin alterar visiblemente el resto del grafo. Esa última frase es el gancho hacia
el ejemplo de la siguiente diapositiva.

% ============================================================
\section{Ejemplo motivador}
\tiempo{1:00}
Leer la pregunta en voz alta, luego señalar la columna izquierda: ese es el camino de
razonamiento legítimo hacia la respuesta correcta (Estados Unidos), tres saltos completos.
Luego la columna derecha: con solo \textbf{una arista falsa} insertada
(\texttt{Manchester --cityOf--> England}), aprovechando una tripleta ya existente y
verídica pero ajena a la película (\texttt{England --containedIn--> United Kingdom}), se
completa una ruta alternativa de dos saltos hacia una respuesta incorrecta. Enfatizar: la
tripleta insertada por sí sola no basta, se apoya en algo que ya estaba en el grafo --- así
es como el ataque se mantiene sigiloso. Cerrar con el objetivo declarado: documentar una
implementación propia de este mecanismo, de extremo a extremo, sobre un caso verificable a
mano.

% ============================================================
\section{Grafo de conocimiento y caminos de relaciones}
\tiempo{0:45}
Tres definiciones que se van a usar sin volver a explicar durante el resto de la charla,
así que conviene ir despacio aquí:
\begin{itemize}
  \item \textbf{Grafo de conocimiento} $G=(E,R,T)$: un conjunto de tripletas
    $(e_h, r, e_t)$, cada una una arista dirigida y etiquetada de una entidad cabeza a una
    entidad cola.
  \item \textbf{Camino de razonamiento} $p$: una secuencia de saltos entidad-relación-entidad;
    su \textbf{camino de relaciones} asociado $w=(r_1,\dots,r_l)$ es la misma secuencia pero
    descartando las entidades intermedias --- solo importa el patrón de relaciones, no por
    dónde pasa exactamente.
  \item \textbf{Anclaje (\emph{grounding})}: la operación inversa. Dado un camino de
    relaciones $w$ y una entidad de inicio, encontrar qué entidades reales del grafo se
    alcanzan al recorrerlo.
\end{itemize}
El anclaje es la pieza que se vuelve a usar literalmente en la Ecuación 7 (diapositiva 10) y
en la Etapa 3 del ataque (diapositiva 14), así que vale la pena remarcarlo.

% ============================================================
\section{Pipeline de un sistema KG-RAG}
\tiempo{1:00}
Dos ecuaciones, una por etapa de un sistema KG-RAG normal (sin ataque todavía):
\textbf{Ec. 1} ($G_q = \text{Retriever}(q; G)$) es la etapa de recuperación: dado el grafo
completo $G$ y la pregunta $q$, se extrae un subgrafo $G_q$ con solo la información
relevante. \textbf{Ec. 2} ($a = f_\theta(\text{Prompt}[q, G_q])$) es la etapa de
generación: un LLM con parámetros $\theta$ recibe la pregunta y el subgrafo recuperado
combinados en un prompt, y produce la respuesta final $a$. La vulnerabilidad estructural
que sostiene todo el artículo aparece justo aquí: si $G$ contiene tripletas falsas bien
conectadas, la Ec. 1 puede incluirlas en $G_q$ sin ningún mecanismo que las distinga de las
legítimas --- el atacante no necesita tocar la Ec. 2 en absoluto.

% ============================================================
\section{Modelo de amenaza (caja negra)}
\tiempo{1:00}
La Ecuación 3 formaliza el objetivo del atacante:
$\hat{A} = \text{KG-RAG}(q; \hat{G})$, con $a^{*} \notin \hat{A}$, sobre un grafo envenenado
$\hat{G} = \{E, R, T \cup \hat{T}\}$. En palabras: el atacante inserta un conjunto de
tripletas adversarias $\hat{T}$ en el grafo original, de modo que cuando el sistema
KG-RAG completo opera sobre el grafo envenenado $\hat{G}$, el conjunto de respuestas que
produce $\hat{A}$ \textbf{ya no contiene} la respuesta correcta original $a^{*}$. Tres
restricciones importantes de remarcar en voz alta:
\begin{enumerate}
  \item El atacante \textbf{no conoce} el recuperador, el LLM (ni su arquitectura ni sus
    parámetros), ni la respuesta correcta de la pregunta que está atacando --- de ahí
    ``caja negra''.
  \item Su única capacidad es \textbf{insertar} tripletas; no puede borrar ni modificar las
    existentes.
  \item Restricciones de sigilo: cada tripleta insertada reutiliza entidades y relaciones
    \emph{ya existentes} en el grafo (nunca inventa vocabulario nuevo), y el número de
    tripletas por respuesta adversaria está acotado por un presupuesto $K$.
\end{enumerate}

% ============================================================
\section{Formalización: objetivo de extracción de caminos (I--II) --- Ecuaciones 4, 5 y 6}
\tiempo{1:45}
\textbf{Esta es la diapositiva que fusiona dos slides originales (Ec. 4--5 y Ec. 6) porque
las tres ecuaciones son un solo argumento en tres pasos.} Conviene narrarlas como una
cadena, no como tres fórmulas sueltas:

\begin{description}[leftmargin=1.2em]
  \item[Ecuación 4 --- el objetivo ideal.]
    $\displaystyle\max_\theta \; \mathbb{E}_{w \sim W^{*}}\big[\log P_\theta(w \mid q)\big]$.
    Se busca el conjunto de parámetros $\theta$ que maximiza la probabilidad esperada, bajo
    la distribución de caminos ``correctos'' $W^{*}$, de que el modelo genere ese camino de
    relaciones $w$ dada la pregunta $q$. $W^{*}$ es el conjunto de caminos de relaciones
    \textbf{más cortos} entre la entidad tema y la respuesta correcta en el grafo --- se le
    llama la \emph{supervisión débil} del problema porque no hay etiquetas manuales de
    ``el camino correcto es este''; se deriva algorítmicamente calculando caminos mínimos.
  \item[Ecuación 5 --- la versión computable.]
    $\displaystyle\operatorname*{arg\,max}_\theta \; \frac{1}{|W^{*}|} \sum_{w \in W^{*}} \log P_\theta(w \mid q)$.
    La Ec. 4 tiene una esperanza matemática que no se puede calcular directamente sin
    conocer la distribución completa de $w$. La Ec. 5 la resuelve asumiendo que $w$ es
    uniforme dentro de $W^{*}$: bajo ese supuesto, la esperanza se aproxima simplemente
    \textbf{promediando} el log-verosímil sobre todos los caminos del conjunto finito
    $W^{*}$. Es literalmente la función de pérdida que se escribiría en código --- el paso
    de ``qué queremos'' (Ec. 4) a ``qué se puede entrenar'' (Ec. 5).
  \item[Ecuación 6 --- la descomposición práctica.]
    $\displaystyle\frac{1}{|W^{*}|} \sum_{w \in W^{*}} \log \prod_{i=1}^{|w|} P_\theta(r_i \mid r_{<i}, q)$.
    Un camino de relaciones $w=(r_1,\dots,r_l)$ es una \textbf{secuencia}, así que su
    probabilidad $P_\theta(w\mid q)$ de la Ec. 5 se factoriza con la regla de la cadena en
    el producto de las probabilidades de cada relación individual $r_i$, condicionada a
    todas las relaciones generadas antes ($r_{<i}$) y a la pregunta $q$. Esto convierte el
    problema en algo idéntico a modelado de lenguaje estándar, pero prediciendo
    \emph{relaciones del grafo} en vez de palabras: entrenar el modelo para proponer buenos
    caminos equivale a entrenarlo para predecir, relación por relación, la siguiente más
    plausible.
\end{description}

Cerrar con el puente hacia las desviaciones: el artículo resuelve este objetivo
\textbf{entrenando} un LLM específico de KG (LLM\_RoG, basado en LLaMA-2-7B) con estas tres
ecuaciones como función de pérdida. Esta implementación logra el mismo efecto --- caminos
anclables en el grafo real --- sin entrenar nada, restringiendo el vocabulario de un LLM de
propósito general (detalle completo en la diapositiva 18).

% ============================================================
\section{Formalización: tripleta de perturbación --- Ecuación 7}
\tiempo{1:00}
$\displaystyle E_{l-1} = \text{Grounding}(e_q, w'; G)$, con $w' = (r_1, \dots, r_{l-1})$.
Aquí se retoma el anclaje definido en la diapositiva 6. $w'$ es el \textbf{prefijo} del
camino de relaciones $w$: todas sus relaciones menos la última ($r_l$). Se ancla ese
prefijo desde la entidad tema $e_q$, recorriendo
$e_q \xrightarrow{r_1} \cdots \xrightarrow{r_{l-1}} e_{l-1}$ sobre el grafo real, y se
obtiene el conjunto $E_{l-1}$ de entidades efectivamente alcanzadas. Por cada entidad
alcanzada, se construye la tripleta de perturbación
$(e_{l-1}, r_l, \hat{a})$: se adjunta la última relación del camino, apuntando directamente
a la respuesta adversaria $\hat{a}$. En palabras simples: se recorre el camino real hasta el
penúltimo paso, y el último salto se ``clava'' artificialmente hacia la respuesta que se
quiere inducir. \textbf{Estrategia de respaldo}: si el prefijo no se puede anclar en el
grafo (porque ninguna entidad real sigue ese patrón de relaciones), se sintetiza una cadena
con entidades puente aleatorias que preserva el mismo patrón de relaciones hasta llegar a
$\hat{a}$ --- así el ataque nunca se bloquea por falta de anclaje.

% ============================================================
\section{Pipeline de tres etapas}
\recortable
\tiempo{0:40}
Diapositiva puramente de transición: mapea cada etapa del método a su módulo de código
(\texttt{adversarial\_answers.py}, \texttt{relation\_paths.py}, \texttt{perturbation.py}) y
a la función orquestadora \texttt{run\_attack(kg, question, topic\_entity, llm,
n\_answers=5, budget\_k=4)}. Si el tiempo aprieta, se puede pasar directo a la Etapa 1
mencionando de paso que las tres etapas están encadenadas por \texttt{run\_attack} con
$N=5$ y $K=4$ como valores por defecto (los mismos del artículo).

% ============================================================
\section{Etapa 1 --- Generación de respuestas adversarias}
\tiempo{1:00}
Mostrar el prompt (tomado literalmente del artículo): se le pide al LLM cinco nombres de
entidad que sean respuestas \emph{incorrectas} pero plausibles o confusas para la pregunta.
El problema es que el LLM puede alucinar nombres que ni siquiera existen en el grafo actual
--- una respuesta adversaria inútil si la entidad no está ahí para anclar un camino hacia
ella. Por eso se repiten rondas de generación y se aplica \textbf{coincidencia difusa}
(\texttt{difflib.get\_close\_matches}) contra el vocabulario real de entidades del grafo,
hasta reunir $N$ respuestas que sí existen como nodos.

% ============================================================
\section{Etapa 2 --- Extracción de caminos de relaciones}
\tiempo{1:00}
Aquí vive la desviación principal frente al artículo, aunque se detalla recién en la
diapositiva 18 --- en esta diapositiva basta con describir el mecanismo: se calcula el
vocabulario de relaciones alcanzables en pocos saltos desde la entidad tema
(\texttt{neighborhood\_relations}), y se le pide al LLM caminos de relaciones usando
\textbf{solo} ese vocabulario, con un prompt propio de esta implementación (el artículo
resuelve esta etapa entrenando un modelo, no con un prompt). Cualquier camino que use una
relación fuera del vocabulario permitido se descarta en el post-procesamiento. La razón de
que esto funcione: restringir el vocabulario garantiza caminos anclables en el grafo real,
sin necesidad de entrenar un modelo específico de KG.

% ============================================================
\section{Etapa 3 --- Inserción de tripletas de perturbación}
\tiempo{1:00}
Esta etapa es la aplicación directa de la Ecuación 7 (diapositiva 10): estrategia principal
--- ancla el prefijo del camino desde la entidad tema y adjunta la última relación hacia la
respuesta adversaria; estrategia de respaldo --- si el prefijo no ancla, muestrea entidades
puente al azar preservando el mismo patrón de relaciones. Cerrar con las dos reglas
operativas: presupuesto $K$ tripletas por respuesta adversaria, y nunca se duplica una
tripleta que ya exista en el grafo (eso delataría el ataque).

% ============================================================
\section{Configuración de la réplica}
\tiempo{0:45}
Transición hacia la demo: instalación y variable de entorno, y los dos modos de ejecución
del script. Explicar el caso de prueba en una frase: grafo construido a mano con 13
tripletas sobre \emph{Manchester By The Sea}, atacando la pregunta de en qué país se filmó.
Distinguir los dos modos: \texttt{--offline} usa \texttt{StaticLLMClient} con respuestas
fijas --- siempre inserta las mismas 2 tripletas hacia \emph{United Kingdom}, es la
referencia reproducible del proyecto; sin el flag, se llama a la API real de
\texttt{gpt-4o-mini} a temperatura 0.7, sin respuestas fijas, así que cada corrida puede
variar --- eso se muestra en la siguiente diapositiva.

% ============================================================
\section{Ejecución con la API real de OpenAI}
\tiempo{1:45}
\textbf{Esta diapositiva fusiona la ejecución real con su interpretación.} Mostrar primero
la salida del script: con el modelo real, la respuesta adversaria que sobrevivió la
coincidencia difusa fue \emph{Florida}, no \emph{United Kingdom} como en el modo offline
--- porque el modelo real no es determinista y propone un conjunto distinto de países en
cada corrida. Explicar por qué \emph{Florida} sobrevivió: existe en este grafo solo por dos
tripletas ajenas a la película (\texttt{Miami, locatedIn, Florida} y
\texttt{Florida, stateOf, United States}), así que la coincidencia es accidental, del mismo
tipo que la de \emph{United Kingdom} en el modo offline. De los dos caminos de relaciones
propuestos, señalar que \texttt{(filmPlace, starring)} es \textbf{anclable pero
semánticamente incoherente}: adjuntar \texttt{starring} (protagonizada por) a una ubicación
como Manchester no tiene sentido real, aunque la restricción de vocabulario de la Etapa 2 lo
admite sin problema --- es un límite honesto de esa desviación. Cerrar con la
interpretación: con solo \textbf{2 tripletas insertadas} sobre 13 originales se crean dos
rutas de razonamiento alternativas hacia \emph{Florida}, sin ninguna señal estructural que
las distinga de la ruta legítima hacia \emph{United States} --- esto confirma el mecanismo
central del artículo, incluso con un camino semánticamente extraño. Y el límite de
reproducibilidad: sin semilla fija para el muestreo del LLM, esta corrida específica no es
reproducible byte a byte (por eso el modo \texttt{--offline} es la referencia canónica del
proyecto).

% ============================================================
\section{Resultados reportados en el artículo original}
\tiempo{1:15}
Aclarar primero que esta tabla \textbf{no es de esta implementación}: son los resultados
del artículo (Tabla 3, WebQSP) que se muestran como contexto de escala, ya que el caso de
Manchester By The Sea no reproduce estas cifras. Con $N=5$ respuestas adversarias y $K=4$
tripletas cada una (máximo 20 tripletas por pregunta), los cuatro sistemas KG-RAG evaluados
sufren caídas sustanciales de F1 frente a una línea base de modificación aleatoria (que
apenas los afecta). El hallazgo clave del artículo, y el que conecta con la Etapa 2 de este
proyecto: la vulnerabilidad principal está en la etapa de \textbf{recuperación} (más del
87\% de las tripletas adversarias son recuperadas); la etapa de generación con el LLM es
más resistente, pero exhibe una respuesta polarizada --- o rechaza el contenido adversario
por completo, o lo adopta dominantemente, sin término medio.

% ============================================================
\section{Desviaciones y limitaciones de alcance}
\tiempo{1:30}
\textbf{Esta diapositiva fusiona ``Desviaciones frente al método original'' con
``Limitaciones de alcance''.} Primero la única desviación deliberada: el artículo entrena
LLM\_RoG (LLaMA-2-7B) con el objetivo de las Ecuaciones 4--6 (diapositiva 9); esta
implementación logra caminos anclables restringiendo el vocabulario de relaciones de un LLM
de propósito general al vecindario real de la entidad tema, sin GPU ni entrenamiento ---
recalcar que el resto del pipeline (respuestas adversarias por coincidencia difusa,
inserción con estrategia de respaldo) sigue el artículo directamente, sin modificaciones de
fondo. Luego las tres limitaciones de alcance, dichas con transparencia y sin sonar como
excusa:
\begin{enumerate}
  \item No se integra ningún sistema KG-RAG objetivo (RoG, GCR, G-retriever, SubgraphRAG):
    la validación se detiene en el grafo envenenado, verificado por inspección directa, sin
    medir métricas de respuesta final.
  \item Sin benchmarks a gran escala: 13 tripletas construidas a mano, frente a los miles de
    preguntas y subgrafos de WebQSP/CWQ.
  \item Modelo de lenguaje distinto (\texttt{gpt-4o-mini} en vez de GPT-4) --- remitir a la
    siguiente diapositiva para el detalle completo.
\end{enumerate}
Cerrar con el trabajo futuro: integrar un recuperador y generador mínimos para medir el
efecto sobre la respuesta final, y ejecutar sobre un subconjunto real de WebQSP/CWQ.

% ============================================================
\section{Modelos de lenguaje evaluados en el artículo}
\tiempo{1:15}
Esta es la Tabla 5 del artículo: robustez del método RoG en WebQSP variando \textbf{solo}
el LLM de la etapa de generación, con recuperación y prompts fijos. Recorrer la tabla de
forma rápida por categoría, no fila por fila: modelos open-source (Llama-2-7B-chat-hf,
Llama-3.1-8B-Instruct) parten más bajo y caen más en términos absolutos; los closed-source
(GPT-3.5-turbo, GPT-4) y el modelo MoE (DeepSeek-V3-0324) rinden mejor tanto limpio como
atacado; y el modelo específico de KG (LLM\_RoG) logra el mejor desempeño limpio pero
también sufre la \textbf{mayor} degradación relativa bajo ataque (46\%) --- hallazgo
contraintuitivo del artículo: saber más sobre la estructura del grafo también lo hace más
vulnerable a su manipulación. Aclarar que para generar las respuestas adversarias del
propio ataque el artículo usa GPT-4, y que GCR/SubgraphRAG usan GPT-3.5-turbo como backend
de evaluación --- ninguno de estos seis modelos coincide con el usado en esta
implementación. Cerrar con el \texttt{alertblock}: esta implementación usa únicamente
\texttt{gpt-4o-mini} en las dos etapas basadas en LLM (respuestas adversarias y caminos de
relaciones). \textbf{Si preguntan por qué}, la respuesta completa está en la sección de
preguntas frecuentes al final de este documento.

% ============================================================
\section{Conclusiones}
\tiempo{0:45}
Cuatro puntos, dichos con seguridad y sin sobre-explicar (ya se dijo todo antes): se
reprodujo el pipeline de tres etapas de forma ejecutable y verificable, con demostración de
extremo a extremo contra la API real; el caso replicado confirma el mecanismo central del
artículo (una perturbación mínima e indistinguible estructuralmente puede desviar una
cadena de inferencia completa); la simplificación principal --- vocabulario restringido en
vez de modelo entrenado --- preserva la propiedad esencial (caminos anclables) pero no
garantiza coherencia semántica, como mostró el camino \texttt{starring} sobre una ubicación
en la ejecución real; y el paso natural siguiente es integrar un sistema KG-RAG objetivo
para medir el impacto sobre la respuesta final.

% ============================================================
\section{Referencias / Gracias}
\tiempo{0:20}
Leer la referencia en formato APA y agradecer. No hay contenido nuevo que explicar.

% ============================================================
\section{Apéndice --- Mapa de módulos}
\respaldo
Solo mostrar si preguntan ``¿dónde vive cada pieza del código?''. Frase de resumen lista
para usar: \emph{``Las tres etapas del método están en \texttt{adversarial\_answers.py},
\texttt{relation\_paths.py} y \texttt{perturbation.py}; \texttt{kg.py} y \texttt{llm.py} son
utilidades compartidas que las tres usan; \texttt{attack.py} las encadena;
\texttt{demo\_manchester.py} es el único punto de entrada ejecutable desde la
terminal.''}

% ============================================================
\section{Apéndice --- Secuencia de ejecución}
\respaldo
Solo mostrar si preguntan ``¿cómo se ejecuta exactamente el script paso a paso?''. Seguir la
numeración del slide: construcción del grafo, instanciación del LLM (real u offline),
llamada a \texttt{run\_attack} (que internamente llama a las tres funciones de las tres
etapas en orden), y finalmente el resumen impreso comparando \texttt{out\_edges} antes y
después del ataque.

% ============================================================
\section*{Preguntas frecuentes anticipadas}
\addcontentsline{toc}{section}{Preguntas frecuentes anticipadas}

\subsection*{¿Por qué no usaste GPT-4 o GPT-3.5-turbo, como el artículo?}
Cuatro capas de respuesta, de la más fuerte a la más técnica:
\begin{enumerate}
  \item \textbf{Vigencia del modelo.} OpenAI retira GPT-3.5-turbo el 23 de octubre de 2026 y
    GPT-4 está en esa misma ola de deprecaciones de 2026. Construir el proyecto sobre
    modelos que la propia OpenAI está apagando durante su desarrollo no es sostenible;
    \texttt{gpt-4o-mini} es el modelo económico vigente y activamente soportado.
  \item \textbf{Costo y volumen de llamadas.} El pipeline hace múltiples rondas de
    generación por etapa (hasta $N=5$ candidatos con reintentos de coincidencia difusa) y se
    ejecuta repetidamente durante desarrollo y pruebas --- un modelo ``mini'' es
    proporcional a ese volumen.
  \item \textbf{La tarea es acotada, no el rol completo evaluado en el artículo.} Aquí el
    LLM no es el generador final de un sistema KG-RAG completo (eso está fuera de alcance);
    cumple dos tareas auxiliares más simples: listas cortas de nombres de entidad y caminos
    de relaciones ya restringidos a un vocabulario filtrado.
  \item \textbf{Ya está declarado como limitación}, no oculto: no se afirma que
    \texttt{gpt-4o-mini} iguale a GPT-4; se reconoce que puede producir candidatos de menor
    diversidad (diapositivas 18 y 19), pero el mecanismo demostrado es una operación
    estructural sobre el grafo, verificable por inspección directa independientemente del
    modelo que generó los candidatos de entrada.
\end{enumerate}

\subsection*{¿Por qué no clonaste y ejecutaste el código del artículo, como pide el formato
del curso?}
Porque el artículo no publica código: se verificó explícitamente en el PDF, en el HTML de
arXiv y en la página de resumen de arXiv, y no hay repositorio ni enlace a implementación
oficial en ninguno de los tres lugares. Ante la ausencia de código que clonar, la única
opción compatible con el objetivo del curso (reproducir el método) es una implementación
propia construida directamente desde la descripción metodológica y las ecuaciones del
artículo --- que es exactamente lo que documenta este trabajo, con sus desviaciones
declaradas de forma explícita en vez de ocultas.

\subsection*{¿Esto reproduce las cifras de degradación del artículo (Tablas 3, 4 y 5)?}
No, y eso está delimitado desde el inicio como parte del alcance: la validación se detiene
en el grafo envenenado, verificado por inspección directa de las tripletas insertadas, sin
integrar ningún sistema KG-RAG objetivo que produzca una respuesta final medible con esas
métricas. Las tablas del artículo que aparecen en la presentación (diapositivas 17 y 19) se
muestran como contexto de escala y como comparación de modelos, no como resultados propios.

\end{document}
